\documentclass{article}
\usepackage[a4paper, total={6.5in, 10in}]{geometry}
\setlength{\parindent}{0pt}
\usepackage{tcolorbox}
\tcbset{colback=white!94!black, colframe=black!60!white, coltitle = white, boxrule=0.8pt, arc=2mm}
\newtcolorbox{boxs}[2][]{title= #2,#1}
\usepackage{datetime2}
\usepackage[british]{babel}
\usepackage{amsfonts}
\usepackage{amsmath}

\title{\textbf{Probability Lecture 2}}
\author{Robert O'Connell}
\date{\today}
\begin{document}
\maketitle
\subsection*{Conditional Probability}
Use new information to revise a model

\subsubsection*{Defintion of Conditional Probability}
\[
P(A|B) = \frac{P(A \cup B)}{P(B)}
\]

\subsubsection*{Properties of Conditional Probabilities}
\begin{itemize}
    \item \(P(A|B) \geq 0\)
    \item \(P(\Omega |B) = 1\)
    \item \(P(B|B) = 1 \)
    \item If \(A \cap C = \emptyset\), then \(P(A \cup C | B) = P(A|B) + P(C|B)\)
\end{itemize}

We can see that conditional probablities satisfy all the standard probability axioms

\subsection*{The Multiplication Rule}

\[
P(A|B) = \frac{P(A \cup B)}{P(B)}
\]
\[
P(A \cup B) = P(B)P(A|B) = P(A)(B|A)
\]
can be extended to \(N\) events

\subsection*{Total Probability Theorem}
We have
\begin{itemize}
    \item Partition of sample space into \(A_1, A_2,\dots, A_n\)
    \item Have \(P(A_i)\) for every \(i\)
    \item Have \(P(B|A_i)\) for every \(i\)
\end{itemize}

Then
\[
P(B) = \sum_{i = 1}^n P(A_i)P(B|A_i)
\]

\subsection*{Bayes' Theorem}
We have
\begin{itemize}
    \item Partition of sample space into \(A_1, A_2,\dots, A_n\)
    \item Have \(P(A_i)\) for every \(i\)
    \item Have \(P(B|A_i)\) for every \(i\)
\end{itemize}

Then
\[
P(A_i|B) = \frac{P(A_i \cap B)}{P(B)} = \frac{P(A_i)P(B|A_i)}{\sum_j P(A_j)P(B|A_j)}
\]

This is a systematic approach for incorporating new evidence
\\

\subsection*{Bayesian Inference}

Initial Beliefs \(P(A_i)\) on possible cause of an observed event B
 
- Model if the world under each \(A_i: P(B|A_i)\) 
\[
A_i \rightarrow B \ \text{using} \ \text{model} \ \& \ P(B|A_i)
\]

- Draw conclusions about causes
\[
B \rightarrow A_i \ \text{using} \ \text{inference} \ \& \ P(A_i|B)
\]

\textbf{Bayesian inference} reasons from observed effects to probable causes. 
Given initial beliefs $P(A_i)$ and forward models $P(B|A_i)$, we observe $B$ 
and compute $P(A_i|B) \propto P(B|A_i)P(A_i)$. \textit{Example:} observing fever ($B$), 
with causes $A_1=\text{flu}$, $A_2=\text{COVID}$, $A_3=\text{cold}$, 
we combine prior probabilities $P(\text{flu})=0.3$, $P(\text{COVID})=0.1$, 
$P(\text{cold})=0.6$ with likelihoods $P(\text{fever}|\text{flu})=0.9$, $P(\text{fever}|\text{COVID})=0.8$, 
$P(\text{fever}|\text{cold})=0.3$ to determine which cause is most likely responsible.

\end{document}